\chapter{Elementos Pré-textuais no \marcaUnDFTeX}

Este capítulo descreve a implementação. Consulte a Parte I para obrigatoriedade, conteúdo e finalidade acadêmica.

\section{Montagem automática}

O arquivo principal chama, em ordem, capa, folha de rosto, verso, errata, folha de aprovação, elementos opcionais, resumos e sumário. Não duplique esses comandos dentro dos arquivos de conteúdo.

\begin{lstlisting}
\undfcover
\pretextual
\undffolhaderosto
\undfversofolhaderostoopcional
\undferrata
\undffolhadeaprovacao
\end{lstlisting}

\section{Resumos}

Preencha \texttt{obrigatorio/resumo.tex}:

\begin{lstlisting}
\undfsetcontent{resumo}{Texto do resumo.}
\undfsetcontent{palavraschave}{termo 1; termo 2; termo 3}
\end{lstlisting}

No arquivo \texttt{obrigatorio/abstract.tex}, informe título, texto e palavras-chave. Para idioma diferente do inglês, altere também o rótulo:

\begin{lstlisting}
\undfsetcontent{abstracttitle}{RESUMEN}
\undfsetcontent{abstract}{Texto em língua estrangeira.}
\undfsetcontent{keywordslabel}{Palabras clave}
\undfsetcontent{keywords}{término 1; término 2}
\end{lstlisting}

O template acrescenta o ponto ao final das palavras-chave. Por isso, informe os termos separados por ponto e vírgula, sem ponto adicional no argumento.

\section{Dados catalográficos}

O arquivo \texttt{opcional/ficha-catalografica.tex} aceita conteúdo textual, imagem ou PDF conforme os comandos públicos do modelo. Utilize apenas o material fornecido pelo fluxo institucional competente.

\begin{lstlisting}
\dadoscatalograficos{Texto fornecido pela unidade responsável.}
% \fichacatalograficaimagem{figuras/ficha.png}
% \fichacatalograficapdf{figuras/ficha.pdf}
\end{lstlisting}

\section{Errata}

\begin{lstlisting}
\errataref{Referência completa do trabalho.}
\erratatexto{%
  \erratatabela{%
    \erratalinha{16}{10}{auto-clavado}{autoclavado}
  }
}
\end{lstlisting}

\section{Elementos opcionais}

Arquivos vazios ou inexistentes não são impressos. O carregamento é feito pelo comando mostrado a seguir. Preencha apenas os elementos usados e não remova a lógica condicional do arquivo principal sem necessidade.

\begin{lstlisting}
\undfinputoptional{opcional/nome-do-arquivo}
\end{lstlisting}

\section{Listas e sumário}

\begin{lstlisting}
\listadeilustracoes
\listadefluxogramas
\listadetabelas
\undfsumario
\end{lstlisting}

O modelo em branco chama apenas o sumário por padrão. Ative listas quando forem adequadas ao trabalho e recompile até que títulos e páginas se estabilizem.
